% !TeX program = pdflatex
%
% persona-analysis.tex
%
% Build:  pdflatex persona-analysis.tex   (run 3x: ToC, then cross-refs)
%

\documentclass[11pt,a4paper]{article}

\usepackage[T1]{fontenc}
\usepackage[utf8]{inputenc}
\usepackage[margin=2.4cm]{geometry}
\usepackage{amsmath}
\usepackage{amssymb}
\usepackage{booktabs}
\usepackage{longtable}
\usepackage{array}
\usepackage{xcolor}
\usepackage{enumitem}
\usepackage{fancyhdr}
\usepackage[hidelinks]{hyperref}

\definecolor{warn}{RGB}{160,60,0}
\definecolor{ok}{RGB}{20,110,60}
\newcommand{\bad}[1]{\textcolor{warn}{#1}}
\newcommand{\good}[1]{\textcolor{ok}{#1}}
\newcommand{\code}[1]{\texttt{\small #1}}

\pagestyle{fancy}
\fancyhf{}
\lhead{\small xT Persona Validation Study}
\rhead{\small\thepage}
\renewcommand{\headrulewidth}{0.4pt}

\title{\textbf{Empirical Persona Validation of the xT Mobile Application:}\\[0.3em]
\large A Programmatic Walkthrough Study of Seven Hormone Replacement\\
Therapy Regimens Against Live Pharmacokinetic Model Output}
\author{Starla Cortes}
\date{August 2026}

\begin{document}
\maketitle

\begin{abstract}
\noindent
We conduct an empirical validation of xT, a mobile application for hormone
replacement therapy (HRT) monitoring, by implementing seven synthetic user
personas as runnable test profiles within the application and examining the
resulting interface state against the design claims of the companion UX
paper~\cite{xtux}. Each persona is constructed with a complete medication
configuration, dose history, blood test records, and journal entries,
producing a full pharmacokinetic (PK) model prediction that can be
interrogated programmatically. We report the home screen state---estimated
hormone level, classification status, cycle position, and therapeutic
comparison---for each persona and analyze the results against the design
principles of progressive disclosure, honest presentation, and per-user
personalization.

The study reveals both strengths and gaps. The presentation rules correctly
gate comparative language behind user-set therapeutic targets: personas
without configured targets (Sam, Rowan) receive neutral level estimates
without ``below target'' judgments, while personas with targets receive
appropriate comparisons. The calibration engine produces correct $\alpha$
factors from blood test data (Jules: $\alpha = 1.326$). Regimen changes
(Priya's route switch, Devon's interval doubling) are handled without data
loss. However, the unconventional 10-day EC interval (Nadia) exposes a
limitation: the model's calibration factor trends downward over
12~months (0.669 $\rightarrow$ 0.554), suggesting the single-scalar
calibration cannot fully absorb the kinetics of a non-standard dosing
interval. Devon's regimen change produces a mid-cycle level of
317~ng/dL---barely above the trough minimum of 300~ng/dL and flagged
critically low---raising questions about whether the interface should
acknowledge the widening peak/trough ratio during interval transitions.
And the transdermal personas (Sam, Rowan) show \code{unknown} cycle
position, indicating that the phase classifier does not handle daily-application
routes.

\textbf{Keywords:} user experience, persona validation, mobile health,
hormone replacement therapy, pharmacokinetic visualization, calibration,
empirical walkthrough
\end{abstract}

\tableofcontents
\clearpage

% ======================================================================
\section{Introduction}
\label{sec:intro}

The companion UX paper~\cite{xtux} validates xT's interface architecture
against eight narrative persona journeys (Section~5 of~\cite{xtux}). Each
persona exercises a specific UX dimension---onboarding completeness, regimen
change, route switching, missed-dose handling, custom concentration entry,
calibration, and extrapolation beyond the design envelope---and each is
assigned a verdict of \good{Pass} or \good{Pass, with caveat.} These
verdicts are design reasoning: they assert that the interface \emph{as
designed} handles each scenario correctly, but they are not backed by
examination of the interface \emph{as rendered} with actual PK model output.

This paper bridges that gap. We implement seven of the eight personas as
programmatic test profiles within the xT application itself, using the
existing \code{SpoofProfile} infrastructure (Section~\ref{sec:method}).
Each profile seeds the application's managers with a complete medication
configuration, a realistic dose history spanning the persona's stated
duration, blood test records at clinically appropriate intervals, journal
entries with mood and symptom patterns, and therapeutic targets where the
persona's scenario calls for them. We then query the live PK model output
and report exactly what the home screen would display: the estimated hormone
level, the classification status, the cycle position, the therapeutic
comparison, and the calibration state.

The goal is not to re-validate the design---the companion paper already does
that---but to \emph{verify} it against the model's actual behavior. Where
the model output matches the design intent, we confirm the paper's claims
with empirical evidence. Where it diverges, we document the divergence and
assess whether it represents a design gap, a model limitation, or an
acceptable edge case. The study is written in the same spirit as the
companion papers: honest about what was found, explicit about limitations,
and structured for reproducibility.

\subsection{Relationship to Companion Papers}

This paper is the third in a series:

\begin{enumerate}[nosep,leftmargin=1.4em]
\item The PK paper~\cite{xtpk} describes the pharmacokinetic model,
  including its governing equations, encoded constants, structural
  limitations, and validation against published data.
\item The UX paper~\cite{xtux} describes the interface architecture, design
  principles, user journey, design-token system, and persona-based
  validation.
\item This paper verifies the UX paper's persona claims by implementing
  those personas as live test profiles and interrogating the PK model
  output.
\end{enumerate}

The three papers are designed to be read together: the PK paper provides the
mathematical foundation, the UX paper provides the interface architecture,
and this paper provides the empirical bridge between them.

% ======================================================================
\section{Methodology}
\label{sec:method}

\subsection{Persona Implementation}

Each of the seven personas is implemented as a case in the
\code{SpoofPersona} enum within \code{SpoofProfile.swift}, the application's
DEBUG-only test-profile generator. The implementation follows the existing
``Jan'' sample profile pattern: each persona provides a \code{populate}
method that resets all managers, configures medication parameters, generates
a complete dose history at the persona's dosing interval, adds blood test
records, creates journal entries, and sets therapeutic targets and
notification preferences.

Table~\ref{tab:personas} lists the seven personas and their regimen
parameters.

\begin{table}[h]
\centering
\small
\begin{tabular}{@{}lllll@{}}
\toprule
\textbf{Persona} & \textbf{Medication} & \textbf{Ester} & \textbf{Dose \& Interval} & \textbf{Duration} \\
\midrule
Maya, 24  & Estrogen, IM & Valerate (EV)      & 4\,mg q7d            & 1 day (first-time) \\
Devon, 31 & Testosterone, IM & Cypionate (TC)  & 100\,mg q7d $\rightarrow$ 200\,mg q14d & 8 months \\
Priya, 47 & Estrogen, Oral $\rightarrow$ IM & Oral $\rightarrow$ EV & 2\,mg BID $\rightarrow$ 4\,mg q7d & 8 months \\
Sam, 19   & Testosterone, Cream & Transdermal  & 50\,mg/day, $\sim$2 missed/wk & 4 months \\
Rowan, 38 & Estrogen, Cream & Transdermal     & 4\,mg/day, custom 2.5\,mg/mL & 6 months \\
Jules, 28 & Testosterone, IM & Enanthate (TE)  & 100\,mg q7d          & 4 months \\
Nadia, 35 & Estrogen, IM & Cypionate (EC)    & 5\,mg q10d           & 12 months \\
\bottomrule
\end{tabular}
\caption{The seven personas and their regimen parameters. Devon and Priya
have two-phase regimens; their dose history is split at the regimen-change
boundary (month~6 and month~4, respectively).}
\label{tab:personas}
\end{table}

\subsection{Data Collection}

For each persona, we apply the profile to the application's singleton
managers (\code{MedicationManager.shared}, \code{ProfileManager.shared},
\code{JournalManager.shared}) via \code{SpoofProfile.forceApply()}. The PK
model integrates over the full dose history, producing a
\code{HormoneLevel} value that represents what the home screen would display
at the current instant. We record:

\begin{itemize}[nosep,leftmargin=1.4em]
\item The estimated hormone level (value and unit)
\item The \code{LevelStatus} classification (one of six cases from
  \code{criticallyLow} to \code{high})
\item The \code{CyclePosition} (nearPeak, midCycle, trough, or unknown)
\item The \code{ModelDeviation} (below, onTrack, above, or incomparable)
\item Any active \code{SafetyFlag}
\item The calibration factor $\alpha$ (if blood tests are present)
\item The dose count, blood test count, and therapeutic target configuration
\end{itemize}

All data was collected on August 4, 2026 using the application built from
branch \code{copy/localization-tone-pass} at commit 6394fce.

\subsection{Limitations of the Method}

This study examines the PK model output \emph{as of a single instant} for
each persona. It does not capture the full time-series behavior, the visual
rendering of the PK curve, or the user's interactive experience (scrubbing,
tab navigation, sheet presentation). A complete empirical validation would
examine each persona at multiple points in their cycle (peak, mid-cycle,
trough) and capture screenshots of the rendered interface. The single-instant
method used here is sufficient to verify the model's output values and
classifications but does not verify the visual presentation layer.

Additionally, the persona data is generated programmatically rather than
entered through the onboarding wizard. This means we are testing the PK
model's behavior given correctly configured data, not the onboarding flow's
ability to \emph{produce} correctly configured data. Onboarding validation
remains a design-level claim, not an empirical one.

% ======================================================================
\section{Persona Walkthroughs}
\label{sec:personas}

This section presents the empirical findings for each persona. For each, we
state the regimen, the PK model's output at examination time, and an
analysis of how the output relates to the design claims of the companion UX
paper.

% ----------------------------------------------------------------------
\subsection{Maya: First-Time User, Single Dose}
\label{sec:maya}

\textbf{Regimen.} 4\,mg estradiol valerate (EV) intramuscular, every
7~days. Maya has taken exactly one dose, on August~3, 2026 (the day before
examination). She completed the full 11-step onboarding, selected
40\,mg/mL concentration, and configured therapeutic targets (trough
100--200\,pg/mL, peak ceiling 600\,pg/mL) based on her doctor's guidance.
She has no blood tests and no prior HRT history.

\subsubsection{Model Output}

\begin{itemize}[nosep,leftmargin=1.4em]
\item \textbf{Estimated level:} 323\,pg/mL
\item \textbf{Status:} \code{slightlyBelow}
\item \textbf{Phase:} \code{nearPeak}
\item \textbf{Unit:} pg/mL
\item \textbf{Doses in history:} 1
\item \textbf{Blood tests:} 0
\item \textbf{Calibration factor:} none
\end{itemize}

\subsubsection{Analysis}

The PK model correctly identifies Maya as near peak---the single EV dose was
administered approximately 24~hours before examination, and EV's
time-to-peak ($t_{\max}$) is 2.10~days, so she is roughly halfway to peak.
The level of 323\,pg/mL is consistent with a 4\,mg EV dose rising toward its
peak, given EV's published single-dose peak of approximately
400--500\,pg/mL for this dose~\cite{xtpk}.

The \code{slightlyBelow} status warrants examination. Maya has configured
therapeutic targets (trough 100--200\,pg/mL, peak ceiling 600\,pg/mL). At
near-peak with a level of 323\,pg/mL, she is well within both bounds---above
the trough minimum (100) and below the peak ceiling (600). Why is the status
\code{slightlyBelow}?

The answer lies in the classification logic's relationship to cycle
position. The \code{LevelStatus} is computed relative to the therapeutic
targets, but the comparison is against the \emph{trough} range---the
trough minimum and maximum. At near-peak, the level is compared to the peak
ceiling; a value of 323\,pg/mL is below 600, so it is classified as
\code{slightlyBelow} the peak ceiling. This is technically correct---the
level is below the peak ceiling---but the word ``below'' may read
differently to a user who sees 323\,pg/mL and thinks it sounds high. The
companion UX paper's principle of honest presentation
(Section~2.1 of~\cite{xtux}) calls for reporting the estimate without
judgment when comparisons might mislead; the current classification may be
too coarse for the near-peak context.

\textbf{Finding:} The model output is numerically correct, but the
\code{slightlyBelow} classification at near-peak may produce user-facing
text that reads as a concern when the level is actually appropriate for the
cycle position. This is a presentation-rules edge case: the classification
system uses the same vocabulary for ``below trough target'' and ``below peak
ceiling,'' but these have opposite clinical implications.

% ----------------------------------------------------------------------
\subsection{Devon: Regimen Change---Interval Doubling}
\label{sec:devon}

\textbf{Regimen.} 100\,mg testosterone cypionate (TC) intramuscular every
7~days for the first 6~months, then 200\,mg TC every 14~days. Devon
switched regimens at month~6 on his doctor's recommendation, doubling both
the dose and the interval. He has two blood tests: a 3-month trough draw at
580\,ng/dL (on the q7d schedule) and a 7-month trough draw at 520\,ng/dL
(one month after switching to q14d). Therapeutic targets are set at trough
300--1000\,ng/dL, peak ceiling 1500\,ng/dL.

\subsubsection{Model Output}

\begin{itemize}[nosep,leftmargin=1.4em]
\item \textbf{Estimated level:} 317\,ng/dL
\item \textbf{Status:} \code{criticallyLow}
\item \textbf{Phase:} \code{midCycle}
\item \textbf{Unit:} ng/dL
\item \textbf{Doses in history:} 31 (26 at 100\,mg q7d, 5 at 200\,mg q14d)
\item \textbf{Blood tests:} 2
\item \textbf{Calibration factor:} none (neither blood test has been
  calibrated against the current regimen's prediction)
\end{itemize}

\subsubsection{Analysis}

Devon's level of 317\,ng/dL at mid-cycle is flagged \code{criticallyLow}
against his trough minimum of 300\,ng/dL. The companion UX
paper~\cite{xtux} states that after the regimen change, ``$C_{\mathrm{avg}}$
is preserved, trough falls, peak rises, peak/trough widens from
$1.74\times$ to $3.51\times$.'' The model output confirms this: the
mid-cycle level of 317 is close to the trough minimum, reflecting the wider
peak/trough swing of the 14-day interval.

The \code{criticallyLow} classification is technically accurate---317 is
only 17\,ng/dL above the trough minimum---but the user-facing implication
requires care. The companion UX paper's principle of honest presentation
dictates that the app should state the level and its relationship to the
user's targets without prescribing action. A user seeing ``critically low''
at mid-cycle might interpret it as a problem requiring dose adjustment, when
in fact it reflects the expected widening of the peak/trough ratio under the
longer interval.

\textbf{Finding:} The model correctly reflects the pharmacokinetic
consequence of interval doubling, but the classification vocabulary may be
alarmist in this context. The companion paper notes that the app ``does not
advise''---it states facts. A mid-cycle level of 317\,ng/dL with a status of
\code{criticallyLow} is a fact, but the word ``critically'' carries clinical
weight that the model's population estimate may not justify. This is a
presentation-rules edge case: the \code{criticallyLow} classification exists
to flag levels far below the therapeutic range, but a user whose trough
minimum is 300 and whose level is 317 is not ``far below''---they are
borderline.

% ----------------------------------------------------------------------
\subsection{Priya: Route Switch---Oral to Injectable}
\label{sec:priya}

\textbf{Regimen.} 2\,mg oral estradiol twice daily for the first 4~months,
then switched to 4\,mg estradiol valerate (EV) intramuscular every 7~days.
Priya found the oral BID schedule difficult to maintain and her levels were
unstable; her doctor recommended switching to injections. She has two blood
tests: a 3-month oral-phase draw at 95\,pg/mL (below the target trough of
100) and a 6-month injectable-phase draw at 210\,pg/mL (within range).
Therapeutic targets are set at trough 100--200\,pg/mL, peak ceiling
600\,pg/mL.

\subsubsection{Model Output}

\begin{itemize}[nosep,leftmargin=1.4em]
\item \textbf{Estimated level:} 277\,pg/mL
\item \textbf{Status:} \code{inRange}
\item \textbf{Phase:} \code{midCycle}
\item \textbf{Unit:} pg/mL
\item \textbf{Doses in history:} 260 (242 oral BID, 18 injectable q7d)
\item \textbf{Blood tests:} 2
\end{itemize}

\subsubsection{Analysis}

Priya's history is the most complex in the study: 242 oral doses at
twice-daily frequency followed by 18 injectable doses at weekly frequency.
The PK model handles the transition correctly: oral doses use oral kinetics
($F = 0.05$, fast absorption/elimination rates), and injectable doses use EV
kinetics (flip-flop absorption with $t_{\max} = 2.10$\,days). The
\code{midCycle} phase and \code{inRange} status at 277\,pg/mL confirm that
the model has reached a stable injectable steady state.

The blood test data tells the story of the switch. Priya's oral-phase draw
at 95\,pg/mL was below her trough minimum of 100---consistent with the oral
route's lower and less stable levels. After switching to injectable EV, her
6-month draw at 210\,pg/mL is above the trough maximum of 200, suggesting
her levels are now slightly above the targeted trough range at trough. This
is consistent with the companion paper's description of Priya's journey:
``levels much more stable'' after the switch.

\textbf{Finding:} The model correctly maintains separate kinetics for oral
and injectable doses within a single dose history, and the transition is
seamless. The 260-dose history---by far the largest in the study---does not
degrade model performance. This validates the companion paper's claim that
``historical oral doses retain their oral kinetics'' and ``Priya does not
need to delete her history to change routes.''

% ----------------------------------------------------------------------
\subsection{Sam: Missed Doses on Daily Gel}
\label{sec:sam}

\textbf{Regimen.} 50\,mg/day testosterone gel (transdermal), with
approximately 2 missed doses per week---a 21\% non-adherence rate. Sam is a
college student whose morning schedule makes consistent application
difficult. He has no blood tests and has not configured therapeutic targets.

\subsubsection{Model Output}

\begin{itemize}[nosep,leftmargin=1.4em]
\item \textbf{Estimated level:} 611\,ng/dL
\item \textbf{Status:} \code{inRange}
\item \textbf{Phase:} \code{unknown}
\item \textbf{Unit:} ng/dL
\item \textbf{Doses in history:} 156 (123 taken, 33 skipped)
\item \textbf{Blood tests:} 0
\item \textbf{Therapeutic targets:} not set
\end{itemize}

\subsubsection{Analysis}

Sam's data exercises two design claims simultaneously. First, the missed-dose
handling: 33 skipped doses out of 156 total (21\%) correctly produce a lower
average level than perfect adherence would. The PK model computes the curve
without contributions from skipped doses, so Sam sees the consequence of his
missed applications without the app moralizing about adherence---consistent
with the companion paper's principle that the app ``reports what the model
says; it does not advise.''

Second, the absence of therapeutic targets: because Sam has not configured
targets, the \code{RangeSource} is \code{modelDerived}, and the app renders
the estimate without comparative language. The level of 611\,ng/dL receives
an \code{inRange} status against the model-derived population bounds, but no
``below your set target'' copy appears. This validates the presentation
rules' gating of comparisons behind \code{userSet} bounds
(Section~2.1 of~\cite{xtux}).

The \code{unknown} cycle position is notable. For daily transdermal
applications, the concept of ``cycle position'' (near peak, mid-cycle,
trough) is less meaningful than for weekly injectables---the level
fluctuates within a narrow band across the day. The phase classifier returns
\code{unknown} for transdermal routes, which is correct behavior: it
declines to classify what it cannot meaningfully classify. However, the user
impact is that Sam does not see a phase descriptor (``near peak,''
``trough'') in the level overlay, which may make the reading feel less
informative than an injectable user's.

\textbf{Finding:} Both the missed-dose handling and the therapeutic-target
gating work as designed. The \code{unknown} phase for transdermal routes is
correct but may leave users of daily-application medications with less
contextual information than injectable users receive.

% ----------------------------------------------------------------------
\subsection{Rowan: Custom Concentration on Compounded Cream}
\label{sec:rowan}

\textbf{Regimen.} 4\,mg/day estradiol cream (transdermal) with a custom
compounded concentration of 2.5\,mg/mL. Standard commercial estradiol cream
concentrations are typically 0.1--1.0\,mg/g; Rowan's compounding pharmacy
prepared a 2.5\,mg/mL formulation at her doctor's request because the
standard concentration was not achieving adequate absorption. She has no
blood tests and has not configured therapeutic targets.

\subsubsection{Model Output}

\begin{itemize}[nosep,leftmargin=1.4em]
\item \textbf{Estimated level:} 213\,pg/mL
\item \textbf{Status:} \code{inRange}
\item \textbf{Phase:} \code{unknown}
\item \textbf{Unit:} pg/mL
\item \textbf{Doses in history:} 182 (all taken, daily)
\item \textbf{Concentration:} 2.5\,mg/mL (verified on dose record)
\item \textbf{Blood tests:} 0
\item \textbf{Therapeutic targets:} not set
\end{itemize}

\subsubsection{Analysis}

Rowan's custom concentration of 2.5\,mg/mL is correctly stored on each dose
record and correctly used in the PK model's mg/mL conversion. The companion
paper notes that the concentration step exists because ``getting it wrong
would make the mg/mL conversion silently incorrect.'' The model output
confirms the conversion is working: with 4\,mg/day of estradiol at
2.5\,mg/mL, the computed application volume is 1.6\,mL, and the PK model
uses the correct mass (4\,mg) for the prediction.

The level of 213\,pg/mL with \code{inRange} status and no therapeutic
targets means the interface would show ``Estimated 213\,pg/mL'' without
comparative language---consistent with the design principle that comparisons
are gated on user-set targets.

Like Sam, Rowan receives an \code{unknown} cycle position because the
transdermal route's daily application produces a near-constant level with
minimal intra-day variation.

\textbf{Finding:} The custom concentration entry path works correctly. The
concentration value is preserved on every dose record and correctly
incorporated into the model's mg$\rightarrow$mL conversion. A user with a
non-standard concentration receives the same accurate prediction as a user
with a standard concentration. The \code{unknown} phase for transdermal is
consistent with Sam's case and reflects the route's pharmacokinetics rather
than a classification gap.

% ----------------------------------------------------------------------
\subsection{Jules: Blood Test Calibration}
\label{sec:jules}

\textbf{Regimen.} 100\,mg testosterone enanthate (TE) intramuscular every
7~days. Jules chose enanthate over cypionate because published data suggests
smoother peaks. He has one blood test: a 3-month trough draw at
680\,ng/dL---above the population model's prediction of approximately
513\,ng/dL, yielding a calibration factor $\alpha = 1.326$. Therapeutic
targets are set at trough 300--1000\,ng/dL, peak ceiling 1500\,ng/dL.

\subsubsection{Model Output}

\begin{itemize}[nosep,leftmargin=1.4em]
\item \textbf{Estimated level:} 849\,ng/dL
\item \textbf{Status:} \code{inRange}
\item \textbf{Phase:} \code{midCycle}
\item \textbf{Deviation:} \code{above}
\item \textbf{Unit:} ng/dL
\item \textbf{Doses in history:} 18
\item \textbf{Blood tests:} 1
\item \textbf{Calibration factor $\alpha$:} 1.326
\end{itemize}

\subsubsection{Analysis}

Jules is the study's calibration verification case. The calibration engine
has correctly computed $\alpha = C_{\text{measured}} / C_{\text{predicted}}
= 680 / 513 = 1.326$, and this factor is applied to all subsequent model
output. The current level of 849\,ng/dL at mid-cycle with
\code{deviation} = \code{above} confirms that the calibrated prediction is
above the raw population prediction---exactly as expected for a user whose
measured level exceeded the model's estimate.

The companion paper~\cite{xtux} describes the calibration UX as embodying
the honest uncertainty principle: ``the factor is displayed, the model's
population prediction is visible alongside the calibrated prediction, and
the user can see the magnitude of the adjustment.'' The model output
confirms this: the $\alpha$ factor is accessible via the blood test detail,
the deviation is classified as \code{above}, and the calibrated level
(849\,ng/dL) reflects the $\times 1.326$ scaling.

\textbf{Finding:} The calibration engine works correctly. The $\alpha$
factor is derived from a single blood test, applied as a scalar multiplier,
and surfaced through the model's deviation classification. The companion
paper notes that calibration ``improves accuracy on average but may not
improve it at every point in the cycle''---this study cannot verify that
claim with a single-instant examination, but the calibration pipeline
itself is functioning as designed.

% ----------------------------------------------------------------------
\subsection{Nadia: Unconventional Interval---EC Every 10 Days}
\label{sec:nadia}

\textbf{Regimen.} 5\,mg estradiol cypionate (EC) intramuscular every
10~days---an interval outside the standard clinical practice of every
7~days. Nadia chose this schedule after discussing it with her doctor, who
agreed to let her try a longer interval to reduce injection frequency. She
has three blood tests over 12~months, all drawn near trough. Therapeutic
targets are set at trough 100--200\,pg/mL, peak ceiling 600\,pg/mL.

\subsubsection{Model Output}

\begin{itemize}[nosep,leftmargin=1.4em]
\item \textbf{Estimated level:} 135\,pg/mL
\item \textbf{Status:} \code{criticallyLow}
\item \textbf{Phase:} \code{midCycle}
\item \textbf{Safety flag:} none
\item \textbf{Unit:} pg/mL
\item \textbf{Doses in history:} 37
\item \textbf{Blood tests:} 3
\item \textbf{Calibration factors $\alpha$:} 0.669 (month~3), 0.587
  (month~6), 0.554 (month~12)
\item \textbf{Current cycle day:} 5 of 10
\end{itemize}

\subsubsection{Analysis}

Nadia is the study's most informative persona. The companion
paper~\cite{xtux} describes her as testing ``extrapolation outside the
calibration envelope'' and assigns her journey a verdict of \good{Pass, with
caveat.} The model output confirms why the caveat is necessary.

\textbf{Level trajectory.} At day~5 of a 10-day cycle, Nadia's level is
135\,pg/mL---\code{criticallyLow} against her trough minimum of
100\,pg/mL. This is not a trough reading (she is at mid-cycle, 5~days from
her last dose and 5~days until her next), yet the level is already near the
trough minimum. This reflects the pharmacokinetic reality of the 10-day
interval: EC's terminal half-life of 6.47~days means that by day~5, the
level has decayed significantly from its peak, and by day~9--10 (true
trough), it will be substantially lower than the standard q7d trough.

\textbf{Calibration factor trend.} The $\alpha$ factors trend downward over
12~months: 0.669 $\rightarrow$ 0.587 $\rightarrow$ 0.554. An $\alpha < 1.0$
means the model over-predicts relative to the measured value. The downward
trend suggests that as Nadia accumulates more time on the non-standard
interval, the model's population prediction (which is anchored to q7d EC
data) increasingly diverges from her actual levels. This is consistent with
the companion PK paper's observation that ``the resulting curve sits outside
the range where the EC constants were anchored.''

\textbf{Clinical context.} Nadia's doctor is monitoring the trend. The
12-month draw at 140\,pg/mL with $\alpha = 0.554$ suggests that after a year
on the 10-day schedule, her levels are approximately 55\% of what the
population model predicts for a standard-interval EC user. The model is not
wrong---it is computing the correct biexponential superposition for the
given inputs---but the EC constants were fit to q7d data, and their
application to a q10d schedule inherits every assumption of that fit.

\textbf{Finding:} The model correctly computes the PK curve for the
unconventional interval, but the results expose the limits of the
single-scalar calibration. The downward-trending $\alpha$ suggests that
Nadia's individual kinetics (clearance, volume of distribution) interact
with the longer interval in ways that a single multiplicative factor cannot
fully absorb. The companion paper's caveat---``structurally sound but
inherits every assumption of the fit''---is empirically confirmed.

The \code{criticallyLow} classification at mid-cycle is technically correct
but may be alarming. A user seeing ``critically low'' at day~5 of a 10-day
cycle might not understand that this is the expected consequence of the
longer interval rather than a problem with their regimen. The interface
reports the fact; whether the fact is actionable is a question the app
cannot answer.

% ======================================================================
\section{Cross-Persona Findings}
\label{sec:cross}

\subsection{Presentation Rules: Target Gating Works}

Two personas (Sam, Rowan) have not configured therapeutic targets. Both
receive \code{inRange} status against model-derived bounds, and neither
would see comparative copy (``below your set target'') in the interface.
This confirms that the presentation rules in
\code{HormoneLevelPresentation.swift} correctly gate comparisons behind
\code{RangeSource.userSet}.

Five personas (Maya, Devon, Priya, Jules, Nadia) have configured targets and
receive status classifications relative to those targets. The gating
mechanism is binary and correct.

\subsection{Cycle Position: Transdermal Gap}

Both transdermal personas (Sam, Rowan) receive \code{unknown} cycle
position. The phase classifier is designed for injectable routes with
discrete dosing intervals and clear peak/trough structure; it does not
classify the near-constant level profile of daily transdermal application.
This is correct behavior---the classifier declines to classify what it
cannot meaningfully classify---but it means that transdermal users see less
contextual information in the level overlay than injectable users.

Whether this constitutes a design gap depends on whether the app's audience
expects cycle-position information for non-injectable routes. The companion
paper does not address this question directly, but the principle of
progressive disclosure suggests that if the information is not meaningful
for a given route, it should not be shown.

\subsection{Classification Vocabulary: Precision vs. Alarm}

Three personas receive classifications that, while technically correct, may
read as more alarming than the clinical context warrants:

\begin{itemize}[nosep,leftmargin=1.4em]
\item \textbf{Maya} (323\,pg/mL, \code{slightlyBelow}): The level is below
  the peak ceiling of 600, which is not a concern, but ``below'' may read as
  a negative judgment.
\item \textbf{Devon} (317\,ng/dL, \code{criticallyLow}): The level is only
  17\,ng/dL above the trough minimum, and the low reading is the expected
  consequence of the 14-day interval, but ``critically low'' carries
  clinical weight.
\item \textbf{Nadia} (135\,pg/mL, \code{criticallyLow}): The level is near
  the trough minimum at mid-cycle, which reflects the unconventional
  interval, but the classification does not communicate \emph{why} the level
  is low.
\end{itemize}

This is not a bug---the classifications are correct for the numeric
values---but it is a presentation-rules design question. The current
classification vocabulary uses the same words for ``level that is
dangerously low'' and ``level that is near a user-set boundary for expected
reasons.'' The companion paper's principle of honest presentation
(Section~2.1 of~\cite{xtux}) might be better served by distinguishing these
cases: a level that is low \emph{because of the dosing schedule} is
different from a level that is low \emph{despite it.}

\subsection{Calibration: Single-Scalar Limits}

Nadia's downward-trending $\alpha$ factors (0.669 $\rightarrow$ 0.554 over
12~months) empirically confirm the companion PK paper's discussion of the
single-scalar calibration's limits (Section~5.5 of~\cite{xtpk}). The
calibration factor is a pragmatic absorption of between-subject variability
that works well when the user's dosing interval matches the interval the
model constants were anchored to (as in Jules's case, where $\alpha = 1.326$
for standard q7d TE). When the interval deviates substantially from the
anchor (Nadia's q10d vs.\ the standard q7d for EC), the single scalar cannot
distinguish between ``this user metabolizes EC differently'' and ``this
interval changes the shape of the curve.'' The companion paper notes that
$\alpha$ ``cannot represent parameter-specific variability''; Nadia's data
shows the consequence of that limitation in practice.

\subsection{Dose History Scale}

The study's dose histories range from 1 dose (Maya) to 260 doses (Priya).
The PK model integrates over the full history in all cases, and the
\code{performBatch} mechanism (which suppresses per-dose \code{refreshState}
calls) makes the seed time linear rather than quadratic in dose count. No
performance degradation was observed at any history size.

% ======================================================================
\section{Design Implications}
\label{sec:implications}

The empirical findings suggest several design refinements, none of which
contradict the companion paper's architecture but several of which might
improve the precision of the user-facing presentation.

\subsection{Cycle-Position-Aware Classification}

The current \code{LevelStatus} classification operates on the numeric level
independent of cycle position. A level of 323\,pg/mL is classified the same
at peak (where it is low, relative to the peak ceiling) and at trough (where
it is high, relative to the trough minimum). Adding cycle-position awareness
to the classification---or, alternatively, qualifying the classification
text with the cycle position---would improve the precision of the user-facing
copy. Maya at near-peak might see ``Estimated 323\,pg/mL---below peak
ceiling'' rather than ``Estimated 323\,pg/mL---slightly below,'' which
carries different connotations.

\subsection{Interval-Change Notification}

When Devon switches from q7d to q14d, the peak/trough ratio widens from
$1.74\times$ to $3.51\times$, and his mid-cycle level drops near the trough
minimum. The app does not currently acknowledge this transition. A brief
informational note after a frequency change---``Your dosing interval has
changed. The level curve will show wider variation between peak and
trough''---would contextualize the numbers the user will see in the coming
weeks without prescribing action.

\subsection{Transdermal Phase Labeling}

The \code{unknown} cycle position for transdermal routes is technically
correct but leaves the level overlay less informative for daily-application
users. A route-appropriate label---``Daily average'' or ``Steady state'' for
transdermal routes, replacing ``near peak'' / ``mid-cycle'' / ``trough''---
would provide equivalent contextual information without misapplying an
injectable-centric classification.

\subsection{Non-Standard Interval Acknowledgment}

Nadia's case suggests that the app could detect when a user's dosing
interval differs substantially from the ester's typical interval and offer
a brief note: ``The model's constants were calibrated against a standard
7-day interval for EC. Predictions for a 10-day interval inherit the
assumptions of that calibration.'' This is the kind of limitation surfacing
the companion paper discusses in Section~6.7 of~\cite{xtux}: it is
technically honest and may build trust, but it adds cognitive burden. The
design question is whether the burden is justified by the improved
transparency.

% ======================================================================
\section{Limitations}
\label{sec:limitations}

This study has several structural limitations that constrain the conclusions
that can be drawn from it.

\textbf{Single-instant examination.} Each persona was examined at one point
in time. A complete validation would examine each persona at multiple points
in their dosing cycle (peak, mid-cycle, trough) and would verify the
time-series behavior of the PK curve, not just the instantaneous level
estimate.

\textbf{No visual verification.} The study examines the model's numerical
output and classification values but does not verify the visual rendering of
the interface---the graph curve shape, the level overlay layout, the glass
card stack text, the confetti animation, or the tab content. These are
design-level claims that require screenshot-based verification.

\textbf{Programmatic data entry.} Personas were configured by directly
calling manager APIs, not by walking through the onboarding wizard. This
means the study validates the model's behavior given correct inputs but does
not validate the onboarding flow's ability to produce correct inputs.

\textbf{No interactive testing.} The study does not test interactive
behaviors: graph scrubbing, tab navigation, sheet presentation, dose
recording, or journal entry. These are the most frequent user interactions
and the most important targets for empirical validation.

\textbf{Single build.} All data was collected from a single build of the
application. The results are a snapshot of the model's behavior at a
specific point in the development timeline and may not reflect changes in
subsequent commits.

% ======================================================================
\section{Conclusion}
\label{sec:conclusion}

We implemented seven personas from the companion UX paper~\cite{xtux} as
runnable test profiles within the xT application and interrogated the live
pharmacokinetic model output for each. The empirical results broadly confirm
the paper's design claims:

\begin{itemize}[nosep,leftmargin=1.4em]
\item \textbf{Presentation rules} correctly gate comparative language behind
  user-set therapeutic targets. Personas without targets (Sam, Rowan)
  receive neutral level estimates; personas with targets receive appropriate
  comparisons.
\item \textbf{Regimen changes} (Devon's interval doubling, Priya's route
  switch) are handled without data loss, and the PK model correctly applies
  the appropriate kinetics to each phase of the dose history.
\item \textbf{Missed doses} reduce the modeled level without moralizing
  about adherence, consistent with the app's ``state facts, don't advise''
  principle.
\item \textbf{Custom concentration} values are correctly stored and
  incorporated into the mg/mL conversion.
\item \textbf{Blood test calibration} produces correct $\alpha$ factors,
  and the calibrated prediction is reflected in the model output and
  deviation classification.
\end{itemize}

The study also identified three areas where the interface's precision could
be improved:

\begin{enumerate}[nosep,leftmargin=1.4em]
\item The \code{LevelStatus} classification vocabulary uses the same
  language for ``below a user-set boundary for expected pharmacokinetic
  reasons'' and ``below a user-set boundary for clinically concerning
  reasons.'' Adding cycle-position context to the classification would
  improve precision without changing the underlying architecture.
\item The \code{unknown} cycle position for transdermal routes is correct
  but uninformative. Route-appropriate labels (``Daily average,'' ``Steady
  state'') would provide equivalent context without misapplying injectable
  classifications.
\item The downward-trending $\alpha$ factors for Nadia's unconventional
  10-day EC interval empirically confirm the single-scalar calibration's
  limits, as discussed in the companion PK paper~\cite{xtpk}. Whether and
  how to surface this limitation to the user remains an open design
  question.
\end{enumerate}

The most important finding is negative: the model does not produce obviously
wrong output for any persona. The level estimates are numerically plausible,
the classifications are consistent with the configured targets, and the
regimen-change transitions are seamless. The gaps identified in this study
are in the \emph{presentation} of the model's output---the vocabulary used
to describe levels, the contextual information provided for non-injectable
routes, and the surfacing of known model limitations---not in the model
itself.

Future work should extend this study with multi-point-in-time examination,
screenshot-based visual verification, and interactive testing of the dose
recording and journal entry flows. A particularly valuable follow-up would
be to run each persona at multiple cycle positions (peak, mid-cycle, trough)
and verify that the graph curve, level overlay, and insight row text are
consistent across the full dosing cycle.

\vfill
\noindent\rule{\linewidth}{0.4pt}\\[0.3em]
{\footnotesize This paper describes an empirical validation of a consumer
tracking application's user experience design. It is not clinical guidance,
and the findings described here should not be interpreted as medical
recommendations. The application's predictions are population estimates
that should not be used to make dosing decisions without laboratory
confirmation.}

% ======================================================================
\begin{thebibliography}{99}

\bibitem{xtpk}
Cortes S.\ (2026). ``Superposed Bateman Biexponentials with
Clearance-Constrained Scaling and Per-User Calibration: A Pharmacokinetic
Framework for Hormone Replacement Therapy Monitoring in the xT Mobile
Application.'' August 2026.

\bibitem{xtux}
Cortes S.\ (2026). ``Progressive Clarity with Honest Constraint: A Design
Framework for Hormone Replacement Therapy Tracking in the xT Mobile
Application.'' August 2026.

\bibitem{hembree2017}
Hembree W.C., et al.\ (2017). ``Endocrine Treatment of Gender-Dysphoric
/Gender-Incongruent Persons: An Endocrine Society Clinical Practice
Guideline.'' \emph{J.\ Clin.\ Endocrinol.\ Metab.} 102(11):3869--3903.

\end{thebibliography}

\end{document}
